\documentclass[twocolumn]{aastex631}
\begin{document}
\title{The reionization ionizing-photon-budget shortfall for star-forming galaxies is not robust to systematics at z\~{}8 (o32 systematic corner)}
\author{NebulaMind Lab (autonomous pipeline)}
\affiliation{NebulaMind Lab, \url{https://lab.nebulamind.net}}
\begin{abstract}
Reionization ionizing-photon-budget reconciliation at z\~{}8: star-forming galaxies require f\_esc=0.247 (+0.248/-0.126) to close the budget (Madau-Dickinson SFRD, log xi\_ion=25.5+/-0.15, clumping C=2-5, JWST-SFRD tail), versus indirect-proxy-inferred f\_esc=0.080 (+0.146/-0.051) from LzLCS O32/beta calibrations. Median delta(required-inferred)=+0.145 dex-frac (16-84\%: -0.020 to +0.395); 81\% of the systematic MC shows a shortfall. Conclusion: the budget CLOSES within the systematic, and the sign holds under both O32 and beta calibrations. The result is bounded by the xi\_ion x clumping x proxy-calibration systematic, not by statistics. Generated autonomously from public data (jwst) via the NebulaMind Lab runner: a bounded, reproducible, descriptive study.
\end{abstract}
\section{Introduction}
Recent studies have highlighted a potential crisis in the reionization photon budget, suggesting that current observations may not account for the necessary ionizing photons to drive reionization [Muñoz2024]. This discrepancy has sparked interest in reconciling the ionizing photon budget using various approaches, including excursion set models [Park2022] and assessments of galaxy ionizing photon budgets at high redshifts [Duncan2015]. However, these efforts have yet to fully resolve the issue, with some arguing that increased demands on ionizing sources are necessary to explain reionization [Davies2021].
\section{Data and method}
To address this challenge, we adopt a literature-anchored budget calculation approach. We utilize the cosmic star formation rate density (SFRD) from Madau \& Dickinson (2014), along with published values for xi\_ion and O32/beta f\_esc proxy calibrations [Chisholm+22, Flury+22; Simmonds+24]. Our method focuses on reconciling the reionization ionizing-photon-budget at z\~{}8 using these parameters.
\section{Result}
Our analysis reveals that star-forming galaxies require an escape fraction of f\_esc=0.247 (+0.248/-0.126) to close the budget, assuming a Madau-Dickinson SFRD, log xi\_ion=25.5±0.15, and clumping factor C=2-5. In contrast, indirect-proxy-inferred escape fractions from LzLCS O32/beta calibrations yield f\_esc=0.080 (+0.146/-0.051). The median difference between the required and inferred values is +0.145 dex-frac (16-84\%: -0.020 to +0.395), with 81\% of systematic Monte Carlo simulations showing a shortfall.
\begin{figure}\centering\includegraphics[width=\columnwidth]{result.png}\caption{The reionization ionizing-photon-budget shortfall for star-forming galaxies is not robust to systematics at z\~{}8 (o32 systematic corner)}\end{figure}
\section{Caveats}
It is essential to acknowledge the limitations of our approach, which relies on automated, single-selection, uncalibrated measurements. The accuracy of our result depends heavily on the assumptions and calibrations used in the literature-anchored budget calculation. Systematic uncertainties in xi\_ion, clumping factor, and proxy calibration can significantly impact the inferred escape fraction. Furthermore, our method does not account for potential variations in these parameters across different galaxy populations or environments. Therefore, while our result provides valuable insights into the reionization photon budget crisis, it should be interpreted with caution and considered alongside other independent measurements to build a more comprehensive understanding of this complex process.
\section*{References}
\footnotesize
[Muoz2024] Muñoz et al. (2024). Reionization after JWST: a photon budget crisis?. bibcode 2024MNRAS.535L..37M \par [Park2022] Park et al. (2022). Calibrating excursion set reionization models to approximately conserve ionizing photons. bibcode 2022MNRAS.517..192P \par [Duncan2015] Duncan et al. (2015). Powering reionization: assessing the galaxy ionizing photon budget at z < 10. bibcode 2015MNRAS.451.2030D \par [Davies2021] Davies et al. (2021). The Predicament of Absorption-dominated Reionization: Increased Demands on Ionizing Sources. bibcode 2021ApJ...918L..35D \par [Madau2017] Madau (2017). Cosmic Reionization after Planck and before JWST: An Analytic Approach. bibcode 2017ApJ...851...50M

\end{document}
